\section{Experimental design}
\label{sec:methods}
\subsection{Foraging environment and controllers}
Agents move in a square arena with walls and regenerating food. Their observations contain 30 values: eight directional sectors for each of food, other agents and walls; current energy; the previous two-dimensional action; and three body traits. A feedforward controller has two hidden layers of 32 tanh units and two outputs, controlling turn and forward speed. Every method uses this actor architecture. Maximum speed, size and sensing range are fixed at 1, 1 and 10, respectively. 

Energy decreases with time, movement and sensing. A food item supplies 40 energy units; an agent starts with 100 and dies when its energy reaches zero. In ecological mode, agents can reproduce above an energy threshold, paying an energy cost for each offspring; they also die at a maximum age. Appendix~\ref{app:implementation} gives the full settings. Episodic competence evaluation uses the same movement, sensing and food dynamics, without reproduction.

Varied training samples arena size, food supply, layout and regrowth from the ranges in Table~\ref{tab:env}. The fixed environment has side length 60, 40 food items, a uniform layout and regrowth probability 0.03. Validation uses 20 sampled training-range environments and four corner cases. We fix thirty test environments before final evaluation: ten within the training ranges (T1), ten smaller, resource-scarce environments with clustered food (T2), and ten with food drift and seasonal relocation (T3). T3 retains the training ranges for the remaining parameters.

\begin{table}[tb]
\centering\small
\caption{Training and fixed test suites. Each test suite contains ten environments. The canonical fixed-training setting is given in the text. Unused layout parameters have no effect in a given environment.}\label{tab:env}
\begin{tabularx}{\linewidth}{lYYY}
\toprule
Parameter & Varied training / T1 & T2 & T3\\\midrule
Arena side $L$ & 40--80 & 30--40 & 40--80\\
Food items $F$ & 20--60 & 10--20 & 20--60\\
Layouts & Uniform, clusters, patches & Clusters & Uniform, clusters, patches\\
Cluster count $k$ & 2--5 & 2--3 & 2--5\\
Cluster scale $\sigma$ & 3--8 & 2--4 & 3--8\\
Patch radius & 4--10 & Not used & 4--10\\
Regrowth probability & 0.01--0.05 & 0.005--0.01 & 0.01--0.05\\
Food drift & 0 & 0 & 0.05--0.3\\
Season period (steps) & None & None & 100--300\\
\bottomrule
\end{tabularx}
\end{table}


\subsection{Training procedures}
\label{sec:procedures}
\paragraph{Ecological neuroevolution.}
OE starts with 32 agents and permits at most 128. Each agent has its own controller, and offspring inherit the parent's weights with Gaussian perturbations. In varied OE, the environmental parameters change every 2,000 world steps while the population continues. OE-fixed retains the canonical environment. During adaptation, either variant continues evolving in the assigned test environment.

Both variants maintain a population floor of 16. After a step below that floor, missing agents are replaced with mutated offspring of uniformly sampled surviving, non-newborn agents; if none are available, agents alive at the start of the step supply the parents. Replacements receive birth energy and a new position. We count natural and floor births separately.

\paragraph{Evolution strategies.}
ES optimizes the shared actor weights using \texttt{Open\_ES} in evosax \citep{salimans2017evolution,lange2023evosax}, with antithetic perturbations and centered-rank fitness shaping. ES-DR samples training environments; ES-fixed uses the canonical environment. Evaluation uses the strategy mean.

Each candidate is evaluated with 16 policy copies in 500-step worlds without births. Its score is food per initial agent, averaged over its configured episodes. Candidates in a generation share environment draws and reset/step random keys; ES-DR redraws the episode environments each generation.

\paragraph{Proximal policy optimization.}
PPO uses a clipped policy-gradient objective \citep{schulman2017proximal}, a separate value network, and a shared policy across agents in 64 parallel episodic worlds. Food acquisition supplies the reward. PPO-DR samples environment parameters at episode resets; PPO-fixed uses the canonical environment. Training actions are stochastic; competence evaluation uses deterministic actions. We adapt PureJaxRL's continuous-action PPO (Appendix~\ref{app:implementation}).

Each world contains 16 agents and restarts at 500 steps or when all agents die. An agent's death terminates its trajectory. Steps beginning with a dead agent are excluded from losses and advantage normalization; when enabled, reward-normalization statistics are collected over living agents.

Table~\ref{tab:procedures} summarizes deployment and transfer.

\begin{table}[tb]
\centering\small
\caption{Training and transfer procedures. Each row has fixed and varied variants, separately tuned. Body traits and actor architecture are shared.}\label{tab:procedures}
\begin{tabularx}{\linewidth}{lYYY}
\toprule
Procedure & Training signal & Competence deployment & State at adaptation start\\\midrule
OE & Energy-dependent survival/reproduction plus floor replacements & 16 sampled population members & Genomes/alive mask retained; world state and energy reset\\
ES & Rank-shaped episodic food score & 16 copies of mean actor & Mean solution retained; fresh strategy/optimizer\\
PPO & Food reward and clipped policy objective & 16 copies of actor & Actor/critic and reward statistics retained; fresh Adam/rollout state\\
\bottomrule
\end{tabularx}
\end{table}


\subsection{Tuning, budgets and evaluation}
We tune each variant separately with eight random-search trials, each using a distinct seed and $5\times10^7$ nominal living-agent steps. We select the configuration with the highest mean zero-shot food acquisition above the random-policy score across the 24 validation environments, then train it with ten final seeds, 5000--5009, at $2\times10^8$ nominal living-agent steps per run. Selection thus uses zero-shot performance, whereas the primary outcome measures performance during adaptation. Fixed-versus-varied comparisons also include differences in selected hyperparameters (Appendix~\ref{app:implementation}).

OE, OE-fixed, PPO-DR, PPO-fixed and ES-DR form the initial comparison. We added ES-fixed after inspecting these results, before its own tuning and final runs.

One living-agent step counts an agent alive at the beginning of a simulated step. Equal agent-step budgets do not equate wall time, floating-point operations or optimizer updates. Updates finish in indivisible units and may cross a checkpoint or stopping target; we analyze both nominal and actual step counts.

Competence is total food consumed divided by the 16 initial agents in a 500-step episode, averaged across eight evaluation episodes. OE samples population members afresh for each episode, without replacement when at least 16 are alive; ES and PPO use 16 policy copies. Episode and reset seeds are shared across methods, although the random-action reference does not share every subsequent world-step random draw with learned policies.

Adaptation starts independently from the final trained snapshot in each test environment and receives $2\times10^7$ nominal living-agent steps. We measure competence before adaptation and at ten subsequent checkpoints. Table~\ref{tab:procedures} and Appendix~\ref{app:implementation} specify the transferred state.

Persistence is evaluated separately, before adaptation. Thirty-two sampled OE agents or policy copies enter ecological mode for 5,000 world steps, with reproduction but no mutation and no population floor. We record extinction and steps survived. 

\subsection{Outcomes and statistical analysis}
\label{sec:statistics}
Let $c_{mre}(t)$ be competence for method $m$, run $r$ and environment $e$ at normalized adaptation checkpoint $t\in\{0,0.1,\ldots,1\}$. We use the prespecified transformation
\begin{equation}
 z_{mre}(t)=\frac{c_{mre}(t)-q_e}{b_e-q_e},\qquad
 A_{mre}=\frac{1}{10}\sum_{j=1}^{10}\frac{z_{mre}(t_{j-1})+z_{mre}(t_j)}{2},
 \label{eq:auc}
\end{equation}
where $q_e$ is the random-policy score and $b_e$ is the largest score across the five initial configurations, all final runs and all eleven checkpoints in that environment. Every environment has $b_e>q_e$. ES-fixed uses the same anchors.

The primary effect is
\begin{equation}
 D=\IQM\{A_{\mathrm{OE},r,e}:e\in T2\cup T3\}
 -\IQM\{A_{\mathrm{PPO\text{-}DR},r,e}:e\in T2\cup T3\},
\end{equation}
where the interquartile mean (IQM) averages the middle half of the pooled run--environment values. We bootstrap complete training runs, keeping their environment outcomes together, with 50,000 draws and statistics seed 4800 \citep{agarwal2021deep,henderson2018deep}. The prespecified margin is $0.05$ on this reference scale. AUC measures performance over the adaptation interval, including starting competence.

The seven initial secondary comparisons use two-sided Mann--Whitney tests on per-run environment means, with Holm adjustment within the family. The ES comparisons form a separate family. The rank tests compare distributions of run means; the effect estimates compare pooled IQMs. Appendix~\ref{app:inference} gives the comparisons and decision rules.

Bootstrap intervals hold the normalization anchors and thirty test environments fixed. They exclude uncertainty from environment and episode sampling, tuning and anchor selection. Because the maximum anchors come from the same runs, interval coverage is uncertain. We examine anchor re-estimation and paired resampling of corresponding seeds in Appendix~\ref{app:diagnostics}.
